\documentclass[xcolor=table,hyperref={pdfpagemode=FullScreen}]{beamer}
% strukt�rizuojame, kad tas pats dokumentas tikt� ir konspektams ir skaidr�ms
% Tod�l \documentclass yra pa�alinti
%xcolor=svgnames,
%xcolor=table
\beamertemplatenavigationsymbolsempty


%\setbeamercolor{section in head/foot}{fg=white,bg=black}






%\defbeamertemplate*{footline}{infolines theme}
%{
%  \leavevmode%
%  \hbox{%
%  \begin{beamercolorbox}[wd=.333333\paperwidth,ht=2.25ex,dp=1ex,center]{author in head/foot}%
%    \usebeamerfont{author in head/foot}\insertshortauthor
%  \end{beamercolorbox}%
%  \begin{beamercolorbox}[wd=.333333\paperwidth,ht=2.25ex,dp=1ex,center]{title in head/foot}%
%    \usebeamerfont{title in head/foot}\insertshorttitle
%  \end{beamercolorbox}%
%  \begin{beamercolorbox}[wd=.333333\paperwidth,ht=2.25ex,dp=1ex,right]{date in head/foot}%
%    \usebeamerfont{date in head/foot}\insertshortdate{}\hspace*{2em}
%    \insertframenumber{} / \inserttotalframenumber\hspace*{2ex}\hspace*{2em}
%  \end{beamercolorbox}%
%\begin{beamercolorbox}[ht=2.25ex,dp=3.75ex]{section in head/foot}
%\insertnavigation{\paperwidth}
%\end{beamercolorbox}%
%}%
%\vskip0pt%
%}

% slide number in right end of slide
\addtobeamertemplate{navigation symbols}{}{%
    \usebeamerfont{footline}%
    \usebeamercolor[fg]{footline}%
    \hspace{1em}%
    \vspace{-1em}%
    \insertframenumber/\inserttotalframenumber
}




\mode<presentation>
\usepackage{pgf}
\pgfdeclareimage[height=1cm]{image}{file}
\setbeamercolor{title}{fg=red!80!black, bg=red!20!white}
\setbeamertemplate{background canvas}[vertical shading][bottom=red!10,top=blue!10]
\useinnertheme[shadow=true]{rounded}
\useoutertheme{shadow}
\usecolortheme{orchid}
\setbeamerfont{block title}{size={}}
%\setbeamercolor{section in head/foot}{fg=red!10,bg=blue!10}
%\setbeamercolor{section in head/foot}{fg=white,bg=black}
%\setbeamertemplate{headline}{%
%   \begin{beamercolorbox}[ht=2.25ex,dp=3.75ex]{section in head/foot}
%        \insertnavigation{\paperwidth}
%    \end{beamercolorbox}%
%}%
%\mode<presentation>


%\makeatletter
%\setbeamertemplate{headline}{%
%    \begin{beamercolorbox}[ht=2.25ex,dp=3.75ex]{section in head/foot}
%        \insertnavigation{\paperwidth}
%    \end{beamercolorbox}%
%}%
%\makeatother





%\setbeamercovered{transparent}
\setbeamercovered{dynamic}




%\usetheme{progressbar}
%(or \usecolortheme{progressbar}, \usefonttheme{progressbar}, \useoutertheme{progressbar}, \useinnertheme{progressbar})


% Tiek skaidr�ms tiek konspektams naudojami paketai

% svarbiausi paketai
\usepackage{amssymb}
\usepackage{amsmath}
\usepackage{graphicx}

% extended math fonts
\usepackage{yhmath}
\usepackage{url}

%\usepackage[authoryear]{natbib}

\bibliographystyle{apalike}
%\bibliographystyle{plainnat}
%\setcitestyle{square,aysep={}}

%\DeclareGraphicsExtensions{ps,eps,pdf}
\graphicspath{{figures/}}
%\usepackage[unicode]{hyperref}

% diagonal lines in tabulars
\usepackage{makecell}
\newcolumntype{x}[1]{>{\centering\arraybackslash}p{#1}}
\usepackage{tikz}
\newcommand\diag[4]{%
  \multicolumn{1}{p{#2}}{\hskip-\tabcolsep
  $\vcenter{\begin{tikzpicture}[baseline=0,anchor=south west,inner sep=#1]
  \path[use as bounding box] (0,0) rectangle (#2+2\tabcolsep,\baselineskip);
  \node[minimum width={#2+2\tabcolsep-\pgflinewidth},
        minimum  height=\baselineskip+\extrarowheight-\pgflinewidth] (box) {};
  \draw[line cap=round] (box.north west) -- (box.south east);
  \node[anchor=south west] at (box.south west) {#3};
  \node[anchor=north east] at (box.north east) {#4};
 \end{tikzpicture}}$\hskip-\tabcolsep}}





%lietuviskos raides

%\usepackage[latin1]{inputenc}
%\usepackage[utf8]{inputenc}
%\usepackage{fourier}
%\usepackage{times}
%\usepackage[T1]{fontenc}
%\usepackage[QX]{fontenc}

%\usepackage[L7x]{fontenc}
%\usepackage{lmodern}

%\usepackage[latin7]{inputenc}
%\usepackage{textcomp}

%\usepackage[lithuanian]{babel}
%\usepackage{hyperref}


%\input /usr/local/texlive/2007/texmf-dist/tex/plain/amsfonts/cyracc.def

%\input /usr/share/texmf/tex/plain/amsfonts/cyracc.def

%\font\tencyr=wncyr10 at 10pt
%\def\cyr{\tencyr\cyracc}
%\font\tencyrbib=wncyr10 at 9pt
%\def\cyrbib{\tencyrbib\cyracc}
%\font\tencyrit=wncyi10 at 10pt
%\def\cyrit{\tencyrit\cyracc}
%\font\tencyritbib=wncyi10 at 9pt
%\def\cyritbib{\tencyritbib\cyracc}
%\font\tencyrbf=wncyb10 at 10pt
%\def\cyrbf{\tencyrbf\cyracc}
%\font\tencyrbfbib=wncyb10 at 10pt
%def\cyrbfbib{\tencyrbfbib\cyracc}

% pataisymai del neaiskaus paketu susipjovimo raidems i kratkaja ir jo \"e

%\def\cyri{\hbox{\accent"24 i}}
%\def\cyre{\hbox{\accent"20 e}}
%\def\cyrI{\hbox{\accent"24 I}}
%\def\cyrE{\hbox{\accent"20 E}}
\newcommand{\ii}{\mathrm{i}} %imaginary unit
\newcommand{\ee}{\mathrm{e}} %exponent
\newcommand{\dd}{\mathrm{d}} %differential
\newcommand{\cl}[2]{\ensuremath{\mathit{Cl}_{#1,#2}}} % clifford algebra notation

\def\d{\!\cdot\!} %narrow dot product
\def\w{\!\wedge\!} %narrow wedge product
\def\D{\cdot} %wide dot product
\def\W{\wedge} %wide wedge product


\newcommand{\e}[1]{\mathbf{e}_{#1}} % Clifford basis vectors in Cl(3,1)
\newcommand{\Ie}[1]{I\negthinspace\mathbf{e}_{#1}}%bivector in Cl(3,1)
\newcommand{\inv}[1]{\overline{#1}} %space inversion
\newcommand{\ba}{\mathbf{a}} %bold a
\newcommand{\bk}{\mathbf{k}}

\def\sR{\mathsf{R}} %sans serif Rotation
\def\sH{\mathsf{H}} %sans serif Hamiltonian
\def\sF{\mathsf{F}} %sans serif F
\def\sG{\mathsf{G}} %sans serif G
\def\cB{\mathcal{B}} %bivector

\newcommand{\bbR}{\ensuremath{\mathbb{R}}}
\newcommand{\bbP}{\ensuremath{\mathbb{P}}}
\newcommand{\bbC}{\ensuremath{\mathbb{C}}}
\newcommand{\bbH}{\ensuremath{\mathbb{H}}}
\newcommand{\bbF}{\ensuremath{\mathbb{F}}}


\def\m#1{\mathsf{#1}} % general multivector
\newcommand{\reverse}[1]{\widetilde{#1}}
\newcommand{\gradeinverse}[1]{\wideparen{#1}}
\newcommand{\cliffordconjugate}[1]{\wideparen{\widetilde{#1}}}

%amsbok stiliaus pakeitimai lietuvi�kai numeracijai
%\usepackage{mydefs}
%sutrumpinimai
%\input mydef.tex






% Nustatymai konspekt� stiliui
\mode<article>
{
 \usepackage{fullpage}
 \usepackage{pgf}
 \usepackage{url}
}




% Nustatymai projektoriaus stiliui


%\setbeamerfont{title}{size=\tiny, parent=structure}
%\beamertemplatenavigationsymbolsempty
%\setbeamerfont{title}{family=\rmfamily}



%\mode<presentation>
%{
%  \setbeamertemplate{background canvas}[vertical shading][bottom=red!10,top=blue!10]
%  \usetheme{Marburg}
%   \usetheme{Warsaw}
%  \usefonttheme[onlysmall]{structurebold}
%  \usecolortheme[named=Teal]{structure}
%}




%{\usebeamercolor{palette primary}}



%\usetheme{Madrid}

%\setbeamercolor{section in head/foot}{fg=white,bg=black}

%\makeatletter
%\setbeamertemplate{headline}{%
%    \begin{beamercolorbox}[ht=2.25ex,dp=3.75ex]{section in head/foot}
%        \insertnavigation{\paperwidth}
%    \end{beamercolorbox}%
%}%
%\makeatother




%Titulinis puslapis
\title{Inverse of  multivector: Beyond p+q=5 threshold}
\author{A.~Acus$^{*}$ an A.~Dargys$^{**}$\\
\normalsize  A.~Go\v{s}tauto 11, LT-01108 Vilnius, Lithuania\\
\normalsize  $^{*}$Institute of Theoretical Physics and Astronomy,\\
\normalsize Vilnius University,\\ Saul{\.e}tekio 3, LT-10257 Vilnius,
Lithuania\\
\normalsize $^{**}$Center for Physical Sciences and Technology,\\ Semiconductor
\normalsize Physics Institute,\\ Saul{\.e}tekio 3, LT-10257 Vilnius,
Lithuania
}



\title[]{Inverse of  multivector: Beyond p+q=5 threshold}

\author[]{\underline{A.~Acus$^{*}$} and A.~Dargys$^{**}$}

\institute{${}^{*}$Institute of Theoretical Physics and Astronomy,
Vilnius University,\\ Saul{\.e}tekio 3, LT-10257 Vilnius,
Lithuania\newline\noindent
$^{**}$Center for Physical Sciences and Technology, Semiconductor
 Physics Institute,\\ Saul{\.e}tekio 3, LT-10257 Vilnius,
Lithuania
}

\date{Spalio 4-6, 2017, Vilnius \newline \vskip 3pt 42${}^\mathrm{oji}$ Lietuvos Nacionalin\.{e} Fizikos konferencija }


\AtBeginSection[]
{
  \begin{frame}<beamer>
    \tableofcontents[currentsection,currentsubsection]
  \end{frame}
}

\colorlet{redshaded}{red!25!bg}
\colorlet{shaded}{black!25!bg}
\colorlet{shadedshaded}{black!10!bg}
\colorlet{blackshaded}{black!40!bg}

\colorlet{darkred}{red!80!black}
\colorlet{darkblue}{blue!80!black}
\colorlet{darkgreen}{green!80!black}

\def\radius{0.96cm}
\def\innerradius{0.85cm}

\def\softness{0.4}
\definecolor{softred}{rgb}{1,\softness,\softness}
\definecolor{softgreen}{rgb}{\softness,1,\softness}
\definecolor{softblue}{rgb}{\softness,\softness,1}

\definecolor{softrg}{rgb}{1,1,\softness}
\definecolor{softrb}{rgb}{1,\softness,1}
\definecolor{softgb}{rgb}{\softness,1,1}


\setbeamertemplate{headline}{%
\leavevmode%
  \hbox{%
    \begin{beamercolorbox}[wd=\paperwidth,ht=2.5ex,dp=1.125ex]{palette quaternary}%
    \insertsectionnavigationhorizontal{\paperwidth}{}{\hskip0pt plus1filll}
    \end{beamercolorbox}%
  }
}

%\useoutertheme[footline=authortitle]{miniframes}


%-----------


\begin{document}

{
\setbeamertemplate{headline}{}
\begin{frame}
  \titlepage
\end{frame}
}
%---------------------------------------------------------------------------------------------
\lecture{}{}
\author[LNFK42, Vilnius 2017]{A.~Acus and A.~Dargys}
%----      0   ------------------------------------

\section[Introduction;]{Introduction}
\begin{frame}
  \frametitle<presentation>{Problem formulation and notations}
\vskip -10pt\begin{columns}[t]
\begin{column}{1.05\textwidth}
\begin{block}{}
\begin{itemize}\small 
\item General Clifford algebra multivector (MV)  
  $\m{A} = \textrm{scalar} + \textrm{bivector} + \cdots + \textrm{pseudoscalar}$
will be considered  in this report 
\item The problem is to find inverse of general MV $\m{A}$, i.~e. $\m{A}^{-1}$ for which 
     $\m{A} \m{A}^{-1} = \m{A}^{-1}\m{A} = 1$ holds,
\item $n = p+q$  in the title indicates that formula for $\m{A}^{-1}$ depends only on dimension of the Clifford algebra $\cl{p}{q}$ and not on particular signature $(p,q)$
\item  The search of solution of the problem heavily relies on computer-aided symbolic mathematics
\end{itemize}
\end{block}
So before explaining calculation of inverse of MV a brief outlook of the {\it Mathematica} package that was written for this purpose will be presented.
\end{column}
\end{columns}
\end{frame}

% ------ 1 ---------
\begin{frame}
  \frametitle<presentation>{GeometricAlgebra package for {\it Mathematica}}
\vskip -20pt
\begin{columns}[t]
\begin{column}{0.6\textwidth}
\begin{block}{Main features}
\begin{itemize}\small 
  \item Textbook notation (implemented without {\it Mathematica} palettes) and automatic precedence of inner, outer and geometric products
\item Algebraic operations in orthonormal frame (additive representation, symbolic coefficients)
\item Switching between multiple algebras (denoted by different color) in same {\it Mathematica} session
\item    Matrix representations of $\cl{p}{q}$ algebras
\item    Idempotents with different base element sorting
\item    Main involutions and general multivector inverse
\end{itemize}
\end{block}
\end{column}
\begin{column}{0.44\textwidth}
\begin{block}{Files}
\begin{itemize} \tiny
  \item GAC\_Base.nb (main file with installation information)
  \item UnicodeCharactersAdd.tr (add-on for {\it Mathematica} system file)
  \item Start.nb (first usage example with short introduction)
  \item AlgebraFormulary.nb (Clifford algebra formulary)
  \item ElementaryFunctions.nb (Clifford function expansions)
  \item Idempotents.nb (algebras idempotents)
  \item Blades (general blade formulas test)
  \item MultivectorInverse (general multivector inverses and formula test)
  \item MatrixRepresentations.nb (Clifford algebra matrix representations)
\end{itemize} 

{\color{black}\small  May be download under licence GPL-3.0 from \url{https://github.com/ArturasAcus/GeometricAlgebra}}
\end{block}
\end{column}
\end{columns}
\end{frame}

%--2----------------
\section[Grade-negated self-product;]{Grade-negated self-product}
\begin{frame}
\frametitle<presentation>{Grade-negated self-product}
\begin{columns}[t]
\begin{column}{0.5\textwidth}
\begin{block}{Grade-$r$ negation operation}
\begin{itemize}\small
\item  Changes the sign of grade-$r$ part of the multivector~$\m{A}$
\item It is an {\color{blue} involution} denoted as $\m{A}_{\bar r}$. Note the bar over negated grade index.
\item Formal definition $\m{A}_{\bar r}=\m{A}-2\langle \m{A}\rangle_r$
\item Main properties: $(\m{A}_{\bar p})_{\bar r}=\m{A}_{{\bar
p},{\bar r}}=\m{A}_{{\bar r},{\bar p}}$, $\m{A}_{{\bar r},{\bar
r}}=\m{A}$, $(\m{A}+\m{B})_{\bar r}=\m{A}_{\bar r}+\m{B}_{\bar
r}$
\item However, $(\m{AB})_{\bar r}\ne \m{A}_{\bar r}\m{B}_{\bar r}$
\item Commutator $\m{A} \m{A}_{\bar
r}-\m{A}_{\bar r}\m{A}=2(\langle \m{A}\rangle_r\m{A}-\m{A}\langle
\m{A}\rangle_r)$
\end{itemize}
\end{block}
\end{column}

\begin{column}{0.5\textwidth}
\begin{block}{Relation to standard involutions}
\begin{itemize}\small
\item Reversion $\reverse{\m{A}}=\m{A}_{{\bar 2},{\bar 3},{\bar 6},{\bar 7},{\bar
{10}},{\bar {11}}\dots}$

\item Grade inversion $\gradeinverse{\m{A}}=\m{A}_{{\bar
1},{\bar 3},{\bar 5},{\bar 7},{\bar {9}},{\bar {11}}\dots}$
\item 
Clifford conjugate $\cliffordconjugate{\m{A}}=\m{A}_{{\bar 1},{\bar 2},{\bar 5},{\bar
6},{\bar {9}},{\bar {10}}\dots}$\\
\end{itemize}


{\color{red}Important: \color{black}  grades $0;4;8;12;\cdots$ are absent in the above list
of standard GA involutions. They play special role in inverse formulas for $n>5$.} 
\end{block}
\end{column}
\end{columns}
 \end{frame}
%------- 3 --------------

\begin{frame}
\frametitle<presentation>{Grade-negated self-product}
\begin{block}{Definitions and the idea}


{\color{black}\small {\color{blue} Grade-negated self-product} is defined via grade-$r$ negation  operation. } 
\begin{itemize}
  \item {\color{black}\small Grade-negated self-product of general multivector $\m{A}$ is geometric product of $\m{A}$ with any number of its grade-negated
      counterparts. Examples ${\color{red}\m{A}} \m{A}_{\bar r} \m{A}_{\bar
      k,\bar l,\bar t} \m{A}_{\bar s, \bar t}\cdots$ {\color{blue}or} $\quad\m{A}_{\bar r}
\m{A}_{\bar k,\bar l,\bar t} \m{A}_{\bar s, \bar t}\cdots{\color{red}\m{A}} $.
operation.}

\item Self-products where the initial MV $\m{A}$ stands in
the left-most or in the right-most position are of special
importance in finding the inverse multivector $\m{A}^{-1}$.

\item Any identity, which contains only operations of geometric product and grade negations, and which was derived in particular Clifford algebra of dimension $n=p+q$ will also be true in any $n$-dimensional algebra of arbitrary signature $(p,q)$. 

\item Assume we can construct a {\color{red} scalar} $s_m$ using just grade-negated self-products of $m$ multivectors with the initial multivector $\m{A}$ standing either in the left-most or right-most position. 
\end{itemize}
\end{block}
\end{frame}



%------- 4 --------------
\begin{frame}
\frametitle<presentation>{Grade-negated self-product}
\begin{block}{Definitions and the idea (continued)}
\begin{itemize}
  \item Name this way obtained scalar the {\color{red}discriminant} of $\m{A}$.
$ {\color{red} \textrm{scalar}:}\quad s_m=\overbrace{\m{A}\m{A}_{\bar r,\bar t,\ldots}\m{A}_{\bar
s,\ldots}\cdots}^m$
\item If we succeed in
constructing such a scalar then left/right inverse of MV can be obtained  by removing left or
right most positioned initial MV $\m{A}$ from the scalar
$s_m=\m{A}\m{A}_{\bar r,\bar t,\ldots}\m{A}_{\bar s,\ldots}\cdots$
and then dividing the result by $s_m$:
\end{itemize}
\begin{equation*}
\color{red}
\m{A}^{-1}=\frac{\m{A}_{\bar r,\bar t,\ldots}\m{A}_{\bar
s,\ldots}\cdots}{s_m}.
\end{equation*}
\end{block}
\end{frame}
%--5--------------------------------
\section[Inverse in $n=p+q\le5$;]{Inverse of general MV of algebras of dimension $n=p+q\le5$}
\begin{frame}
\frametitle<presentation>{Inverses in dimensions $n=p+q\le5$}
\begin{columns}[t]
\begin{column}{0.5\textwidth}
  \begin{block}{Two dimensional algebras: \cl{2}{0} case}\small 
    \begin{itemize}
\item Write $\m{A} = \sum_{J=0}^{2^n-1} a_J\e{J}$, where
the multi-index $J$ covers all orthonormal base elements arranged
in inverse degree lexicographic order, i.~e., the lowest grade
elements appear first while elements of the same grade are ordered
lexicographically: $J=[\{\},\{1\},\{2\},\{1,2\}]$. 

\item Check that self-product $\m{A} \m{A}_{\bar 1,\bar 2}$ immediately gives the required scalar $s_2=a_{\{\}}^2 - a_{1}^2 -
a_{2}^2 + a_{1,2}^2$.
\end{itemize}
%
\end{block}
\end{column}
\begin{column}{0.55\textwidth}
 \begin{block}{Coefficient notation}
\begin{itemize}
\item
It is more convenient to rewrite the coefficient indices as $s_2=a_{0}^2 -
a_{1}^2 - a_{2}^2 + a_{3}^2$, i.~e. the index of scalar $a_{0}$ is zero, the
coefficients at vectors run from $1$ to $n$ while numerical
indices of higher grades increase monotonically up to~$2^n$. 
\item General rule for grade-$r$ index: start the element index by $\sum_{k=0}^{r-1} (\begin{smallmatrix}n\\
k\end{smallmatrix})$ and end by
$\Bigl(\sum_{k=0}^{r-1} (\begin{smallmatrix}n\\
k\end{smallmatrix})\Bigr)+(\begin{smallmatrix}n\\
r\end{smallmatrix})-1$.  Here $(\begin{smallmatrix}n\\
k\end{smallmatrix})$ is the binomial coefficient.
\end{itemize}
\end{block}
\end{column}
\end{columns}
\end{frame}

%---- 6 -----------
\begin{frame}
\frametitle<presentation>{Inverses in dimensions $n=p+q\le5$ }
\vskip -6pt
\begin{columns}[t]
\begin{column}{1.05\textwidth}
\begin{block}{Notes about discriminant}
\begin{itemize}\small
\item <1->The discriminant $s_m$ comprises all coefficients $a_i$ of multivector~$\m{A}$

\item <2->The MV has inverse if discriminant is unequal to zero, $s_m\ne 0$. This imposes
 the requirement for MV coefficients.

\item <3->Because reverse leaves the discriminant invariant,
the formula for $n=2$ can be rewritten as
$s_2=\m{A} \m{A}_{\bar 1,\bar 2}=\m{A}_{\bar 1,\bar 2}\m{A}$, from
which the left/right inverse formulas follow,\vskip 10pt
$\qquad\qquad \m{A}^{-1}= \m{A}_{\bar 1,\bar
2}/(\m{A}\m{A}_{\bar 1,\bar 2});\qquad \qquad
\m{A}^{-1}=\m{A}_{\bar 1,\bar 2}/(\m{A}_{\bar
1,\bar 2}\m{A}).
$\vskip 10pt
\item <4-> The eliminated grades are grade-negated in one of MVs in the self-product.
 In the example the grades $1$ and $2$ have been eliminated simultaneously
 by  $\m{A}_{\bar 1,\bar 2}$. This rule is valid for discriminants  up to $n\le5$.
\item <5->In all cases the highest degree of coefficients in the discriminant
  always matches the degree of determinant of real matrix representation of $\m{A}$, in this case $2\times2$ real matrices ({\color{red} next slide})
\end{itemize}
\end{block}
\end{column}
\end{columns}
\end{frame}
%---- 7 -----------
\begin{frame}
\frametitle<presentation>{Inverses in dimensions $n=p+q\le5$}
\begin{columns}[t]
\begin{column}{1.05\textwidth}
\begin{block}{Period-$8$ table of irreducible
matrix representations of GA for
Cl$_{p,q}$, where $p$ and  $q$ indicate positive and negative signature.}
\begin{center}\label{TableOfPeriodicity1}
%\setlength{\extrarowheight}{1ex}
%\arraycolsep=0pt
%\parbox[t]{0.9\textwidth}{
%\textit{Period-$8$ table of irreducible
%matrix representations of GA (called the fundamental) for
%Cl$_{p,q}$, where $p$ and  $q$ indicate positive and negative signature.}}
%\setlength{\tabcolsep}{1pt}
%\resizebox{\textwidth}{!}{%
%\renewcommand*{\arraystretch}{1}
%\noalign{\vskip -3pt}\hphantom{m}

\resizebox{\columnwidth}{!}{%  
% Acaus preambule
\begin{tabular}{x{0.1em}| >{$}l<{$} >{$}l<{$} >{$}l<{$} >{$}l<{$}
>{$}l<{$} >{$}l<{$} >{$}l<{$} >{$}l<{$}}
\diag{.1em}{.1em}{\mbox{$\!\!p\,$}}{\mbox{\hphantom{n}\smash{$q$}}}
%Dargio preambule
%\begin{tabular}{>{$} l<{$} | >{$} l<{$} >{$}l<{$} >{$}l<{$} >{$}l<{$}
%>{$}l<{$} >{$}l<{$} >{$}l<{$} >{$}l<{$}}
% tabular body
&0&1&2&3&4&5&6&7\\
\cline{2-9}\\[-12pt]
0 &\bbR&\bbC&\cellcolor{red!30}\bbH&{}^2\bbH&\bbH(2)&\bbC(4)&\bbR(8)&{}^2\bbR(8)\\
1 &{}^2\bbR&\cellcolor{red!50}\bbR(2)&\bbC(2)&\bbH(2)&{}^2\bbH(2)&\bbH(4)&\bbC(8)&\bbR(16)\\
2 &\cellcolor{red!50}\bbR(2)&{}^2\bbR(2)&\bbR(4)&\bbC(4)&\bbH(4)&{}^2\bbH(4)&\bbH(8)&\bbC(16)\\
3 &\bbC(2)&\bbR(4)&^2\bbR(4)&\bbR(8)&\bbC(8)&\bbH(8)&^2\bbH(8)&\bbH(16)\\
4 &\bbH(2)&\bbC(4)&\bbR(8)&{}^2\bbR(8)&\bbR(16)&\bbC(16)&\bbH(16)&^2\bbH(16)\\
5 &{}^2\bbH(2)&\bbH(4)&\bbC(8)&\bbR(16)&^2\bbR(16)&\bbR(32)&\bbC(32)&\bbH(32)\\
6 &\bbH(4)&{}^2\bbH(4)&\bbH(8)&\bbC(16)&\bbR(32)&{}^2\bbR(32)&\bbR(64)&\bbC(64)\\
7&\bbC(8)&\bbH(8)&{}^2\bbH(8)&\bbH(16)&\bbC(32)&\bbR(64)&{}^2\bbR(64)&\bbR(128)
\end{tabular}
}
%
\end{center}
\small The left superscript denotes direct sum of matrices,
$^2\bbR(2)\equiv\bbR(2)\oplus\bbR(2)$.
\end{block}
\end{column}
\end{columns}
%\cellcolor{blue!25}
\end{frame}

%---- 8a -----------

\begin{frame}
\vspace{-20pt}
\frametitle<presentation>{Inverses in dimensions $n=p+q\le5$}
\begin{columns}[t]
\begin{column}{1.05\textwidth}
 \begin{block}{Three dimensional algebras: \cl{2}{1} case}
   \begin{itemize} \small 
     \item   Our goal is to eliminate as many grades as possible by forming suitable self-products, for example, for \cl{2}{1} algebra.
      \item The $\m{A}=a_{0}+a_{1} \e{1}+a_{2}
\e{2}+a_{3} \e{3}+a_{4} \e{12}+a_{5} \e{13}+a_{6} \e{23}+a_{7}
\e{123}$ product with $\m{A}_{\bar 1,\bar 2}$ eliminates grades $1$ and $2$.
The result is the new multivector $\m{B}=\m{A} \m{A}_{\bar 1,\bar
2}= b_{0}+b_{7} \e{123}$ which  consists of scalar and grade-$3$
element. In particular, $b_0=a_{0}^2 - a_{1}^2 - a_{2}^2 + a_{3}^2
+ a_{4}^2 - a_{5}^2 - a_{6}^2 + a_{7}^2$ and $b_7=-2 a_{3} a_{4} +
2 a_{2} a_{5} - 2 a_{1} a_{6} + 2 a_{0} a_{7}$.

\item Repeating the grade negation procedure for grade-$3$ then 
yields left-handed discriminant $$s_4=\m{B}
\m{B}_{\bar 3}=\m{A} \m{A}_{\bar 1,\bar 2}(\m{A} \m{A}_{\bar
1,\bar 2})_{\bar 3}=b_0^2 - b_7^2.$$

\item Reversion of $s_4$ immediately gives the $\m{A}$-right-handed discriminant
$$s_4^\prime=(\m{A}_{\bar 1,\bar 2}\m{A})_{\bar 3} \m{A}_{\bar
1,\bar 2}\m{A}$$ that contains four MVs ({\color{red} next slide}).
\end{itemize}
\end{block}
\end{column}
\end{columns}
\end{frame}
%-------- 9 -----------
\begin{frame}
\frametitle<presentation>{Inverses in dimensions $n=p+q\le5$}
\begin{columns}[t]
\begin{column}{1.05\textwidth}
\begin{block}{Period-$8$ table of irreducible
matrix representations of GA for
Cl$_{p,q}$, where $p$ and  $q$ indicate positive and negative signature.}
\begin{center}\label{TableOfPeriodicity1}
%\setlength{\extrarowheight}{1ex}
%\arraycolsep=0pt
%\parbox[t]{0.9\textwidth}{
%\textit{Period-$8$ table of irreducible
%matrix representations of GA (called the fundamental) for
%Cl$_{p,q}$, where $p$ and  $q$ indicate positive and negative signature.}}
%\setlength{\tabcolsep}{1pt}
%\resizebox{\textwidth}{!}{%
%\renewcommand*{\arraystretch}{1}
%\noalign{\vskip -3pt}\hphantom{m}
\resizebox{\columnwidth}{!}{%  
% Acaus preambule
\begin{tabular}{x{0.1em}| >{$}l<{$} >{$}l<{$} >{$}l<{$} >{$}l<{$}
>{$}l<{$} >{$}l<{$} >{$}l<{$} >{$}l<{$}}
\diag{.1em}{.1em}{\mbox{$\!\!p\,$}}{\mbox{\hphantom{n}\smash{$q$}}}
%Dargio preambule
%\begin{tabular}{>{$} l<{$} | >{$} l<{$} >{$}l<{$} >{$}l<{$} >{$}l<{$}
%>{$}l<{$} >{$}l<{$} >{$}l<{$} >{$}l<{$}}
% tabular body
&0&1&2&3&4&5&6&7\\
\cline{2-9}\\[-12pt]
0 &\bbR&\bbC&\bbH&\cellcolor{red!30}{}^2\bbH&\bbH(2)&\bbC(4)&\bbR(8)&{}^2\bbR(8)\\
1 &{}^2\bbR&\bbR(2)&\cellcolor{red!30}\bbC(2)&\bbH(2)&{}^2\bbH(2)&\bbH(4)&\bbC(8)&\bbR(16)\\
2 &\bbR(2)&\cellcolor{red!50}{}^2\bbR(2)&\bbR(4)&\bbC(4)&\bbH(4)&{}^2\bbH(4)&\bbH(8)&\bbC(16)\\
3 &\cellcolor{red!30}\bbC(2)&\bbR(4)&^2\bbR(4)&\bbR(8)&\bbC(8)&\bbH(8)&^2\bbH(8)&\bbH(16)\\
4 &\bbH(2)&\bbC(4)&\bbR(8)&{}^2\bbR(8)&\bbR(16)&\bbC(16)&\bbH(16)&^2\bbH(16)\\
5 &{}^2\bbH(2)&\bbH(4)&\bbC(8)&\bbR(16)&^2\bbR(16)&\bbR(32)&\bbC(32)&\bbH(32)\\
6 &\bbH(4)&{}^2\bbH(4)&\bbH(8)&\bbC(16)&\bbR(32)&{}^2\bbR(32)&\bbR(64)&\bbC(64)\\
7&\bbC(8)&\bbH(8)&{}^2\bbH(8)&\bbH(16)&\bbC(32)&\bbR(64)&{}^2\bbR(64)&\bbR(128)
\end{tabular}
}
%
\end{center}
\small The left superscript denotes direct sum of matrices,
$^2\bbR(2)\equiv\bbR(2)\oplus\bbR(2)$.
\end{block}
\end{column}
\end{columns}
%\cellcolor{blue!25}
\end{frame}
%---- 8b -----------

\begin{frame}
\vspace{-20pt}
\frametitle<presentation>{Inverses in dimensions $n=p+q\le5$}
\begin{columns}[t]
\begin{column}{1.05\textwidth}
  \begin{block}{Three dimensional algebras: \cl{2}{1} case (continued)}
\begin{itemize} \small 
  \item For other 3D algebras, namely \cl{3}{0}, \cl{1}{2}, and \cl{0}{3} the same formula gives different fully expanded expressions, with other signs of individual coefficients~$a_i$.

\item The discriminant is determined by squares of two coefficients, $b_0^2$ and $b_7^2$, the difference of which should not be zero for invertible multivector $\m{A}$ to exist.
\end{itemize}
\end{block}
\end{column}
\end{columns}
\end{frame}
%--10-------------------------------------------------
\begin{frame}
\vspace{-20pt}
\frametitle<presentation>{Inverses in dimensions $n=p+q\le5$}
\begin{columns}[t]
\begin{column}{1.05\textwidth}
 \begin{block}{Four dimensional algebras: $n=p+q=4$}
\begin{itemize} \small 
 \item   In this case elimination of grades can proceed in two alternative ways.
 \item 1) Negate simultaneously grades $1$ and $2$ and then in the next step the grades $3$ and $4$, i.~e. $s_4= {\color{red} \m{A}}\m{A}_{\bar{1},\bar{2}}(\m{A}
\m{A}_{\bar{1},\bar{2}})_{\bar{3},\bar{4}}$
 \item 2) Eliminate $2$ and $3$, and then $1$ and $4$ grades, i.~e. $s_4= {\color{red} \m{A}}
\m{A}_{\bar{2},\bar{3}}
(\m{A}\m{A}_{\bar{2},\bar{3}})_{\bar{1},\bar{4}}$.

\item Similar elimination can be done in right-hand versions,
$s_4^\prime=(\m{A}_{\bar{2},\bar{3}}\m{A})_{\bar{1},\bar{4}}\m{A}_{\bar{2},\bar{3}}
{\color{red} \m{A}}=
(\m{A}_{\bar{1},\bar{2}}\m{A})_{\bar{3},\bar{4}}\,\m{A}_{\bar{1},\bar{2}}
{\color{red} \m{A}}$.

\item Despite different forms, expanded explicit
expressions for discriminants $s_m$ and $s_m^\prime$ in fact always appear to be equal,
$s_m=s_m^\prime$.

\item The appearance of symbol $\m{A}$ as much as four times in
the products $s_4$ and  $s_4^{\prime}$ is again what we expect
from the matrix representation for $n=4$ algebras, ({\color{red} next slide}) namely, the
highest degree of determinant polynomial should be of degree $4$
as well, in accord with $4\times 4$ matrix representation
$\mathbb{R}(4)$. 
\end{itemize}
\end{block}
\end{column}
\end{columns}
\end{frame}
%-------- 11 -----------
\begin{frame}
\frametitle<presentation>{Inverses in dimensions $n=p+q\le5$}
\begin{columns}[t]
\begin{column}{1.05\textwidth}
\begin{block}{Period-$8$ table of irreducible
matrix representations of GA for
Cl$_{p,q}$, where $p$ and  $q$ indicate positive and negative signature.}
\begin{center}\label{TableOfPeriodicity1}
%\setlength{\extrarowheight}{1ex}
%\arraycolsep=0pt
%\parbox[t]{0.9\textwidth}{
%\textit{Period-$8$ table of irreducible
%matrix representations of GA (called the fundamental) for
%Cl$_{p,q}$, where $p$ and  $q$ indicate positive and negative signature.}}
%\setlength{\tabcolsep}{1pt}
%\resizebox{\textwidth}{!}{%
%\renewcommand*{\arraystretch}{1}
%\noalign{\vskip -3pt}\hphantom{m}

\resizebox{\columnwidth}{!}{%  
% Acaus preambule
\begin{tabular}{x{0.1em}| >{$}l<{$} >{$}l<{$} >{$}l<{$} >{$}l<{$}
>{$}l<{$} >{$}l<{$} >{$}l<{$} >{$}l<{$}}
\diag{.1em}{.1em}{\mbox{$\!\!p\,$}}{\mbox{\hphantom{n}\smash{$q$}}}
%Dargio preambule
%\begin{tabular}{>{$} l<{$} | >{$} l<{$} >{$}l<{$} >{$}l<{$} >{$}l<{$}
%>{$}l<{$} >{$}l<{$} >{$}l<{$} >{$}l<{$}}
% tabular body
&0&1&2&3&4&5&6&7\\
\cline{2-9}\\[-12pt]
0 &\bbR&\bbC&\bbH&{}^2\bbH&\cellcolor{red!30}\bbH(2)&\bbC(4)&\bbR(8)&{}^2\bbR(8)\\
1 &{}^2\bbR&\bbR(2)&\bbC(2)&\cellcolor{red!30}\bbH(2)&{}^2\bbH(2)&\bbH(4)&\bbC(8)&\bbR(16)\\
2 &\bbR(2)&{}^2\bbR(2)&\cellcolor{red!50}\bbR(4)&\bbC(4)&\bbH(4)&{}^2\bbH(4)&\bbH(8)&\bbC(16)\\
3 &\bbC(2)&\cellcolor{red!50}\bbR(4)&^2\bbR(4)&\bbR(8)&\bbC(8)&\bbH(8)&^2\bbH(8)&\bbH(16)\\
4 &\cellcolor{red!30}\bbH(2)&\bbC(4)&\bbR(8)&{}^2\bbR(8)&\bbR(16)&\bbC(16)&\bbH(16)&^2\bbH(16)\\
5 &{}^2\bbH(2)&\bbH(4)&\bbC(8)&\bbR(16)&^2\bbR(16)&\bbR(32)&\bbC(32)&\bbH(32)\\
6 &\bbH(4)&{}^2\bbH(4)&\bbH(8)&\bbC(16)&\bbR(32)&{}^2\bbR(32)&\bbR(64)&\bbC(64)\\
7&\bbC(8)&\bbH(8)&{}^2\bbH(8)&\bbH(16)&\bbC(32)&\bbR(64)&{}^2\bbR(64)&\bbR(128)
\end{tabular}
}
%
\end{center}
\small The left superscript denotes direct sum of matrices,
$^2\bbR(2)\equiv\bbR(2)\oplus\bbR(2)$.
\end{block}
\end{column}
\end{columns}
%\cellcolor{blue!25}
\end{frame}
%--12---------------------------
\begin{frame}
\vspace{-20pt}
\frametitle<presentation>{Inverses in dimensions $n=p+q\le5$}
\begin{columns}[t]
\begin{column}{1.05\textwidth}
 \begin{block}{Five dimensional algebras: $n=p+q=5$}
\begin{itemize} \small 
 \item   1) Eliminate simultaneously grades $2$ and $3$ 
 \item 2) Then eliminate grades $1$ and $4$
 \item 3) Finally eliminate grade $5$. Apart from
the last step this sequence is exactly the same as the second
alternative of $n=4$ case. 
\item The final result is
\vskip 10pt
$\qquad\qquad 
s_8^\prime=(\m{F}\m{A})_{\bar{5}}\m{F}{\color{red}\m{A}}\quad $ with
$\quad \m{F}=(\m{A}_{\bar{2},\bar{3}}\m{A})_{\bar{1},\bar{4}}
\m{A}_{\bar{2},\bar{3}}$\vskip 5pt
for right-hand sided form of the
discriminant and 
\vskip 10pt
$\qquad\qquad s_8={\color{red}\m{A}}\m{D}\, (\m{A}\m{D})_{\bar{5}}\quad $ with
$\m{D}=\m{A}_{\bar{2},\bar{3}}
(\m{A}\m{A}_{\bar{2},\bar{3}})_{\bar{1},\bar{4}}$\vskip 5pt
for left-hand
discriminant.

\item Matrix representation of five dimensional Clifford
algebras corresponds to real ${^2\mathbb{R}(4)}$ matrix ({\color{red} next slide}) the the
determinant of which is a polynomial of degree $8$. This is in
accord with the structure of discriminant~$s_8=s_8^\prime$ where
the $\m{A}$ occurs as much as eight times.
\end{itemize}
\end{block}
\end{column}
\end{columns}
\end{frame}
%---------- 13 -------
\begin{frame}
\frametitle<presentation>{Inverses in dimensions $n=p+q\le5$}
\begin{columns}[t]
\begin{column}{1.05\textwidth}
\begin{block}{Period-$8$ table of irreducible
matrix representations of GA for
Cl$_{p,q}$, where $p$ and  $q$ indicate positive and negative signature.}
\begin{center}\label{TableOfPeriodicity1}
%\setlength{\extrarowheight}{1ex}
%\arraycolsep=0pt
%\parbox[t]{0.9\textwidth}{
%\textit{Period-$8$ table of irreducible
%matrix representations of GA (called the fundamental) for
%Cl$_{p,q}$, where $p$ and  $q$ indicate positive and negative signature.}}
%\setlength{\tabcolsep}{1pt}
%\resizebox{\textwidth}{!}{%
%\renewcommand*{\arraystretch}{1}
%\noalign{\vskip -3pt}\hphantom{m}

\resizebox{\columnwidth}{!}{%  
% Acaus preambule
\begin{tabular}{x{0.1em}| >{$}l<{$} >{$}l<{$} >{$}l<{$} >{$}l<{$}
>{$}l<{$} >{$}l<{$} >{$}l<{$} >{$}l<{$}}
\diag{.1em}{.1em}{\mbox{$\!\!p\,$}}{\mbox{\hphantom{n}\smash{$q$}}}
%Dargio preambule
%\begin{tabular}{>{$} l<{$} | >{$} l<{$} >{$}l<{$} >{$}l<{$} >{$}l<{$}
%>{$}l<{$} >{$}l<{$} >{$}l<{$} >{$}l<{$}}
% tabular body
&0&1&2&3&4&5&6&7\\
\cline{2-9}\\[-12pt]
0 &\bbR&\bbC&\bbH&{}^2\bbH&\bbH(2)&\cellcolor{red!30}\bbC(4)&\bbR(8)&{}^2\bbR(8)\\
1 &{}^2\bbR&\bbR(2)&\bbC(2)&\bbH(2)&\cellcolor{red!30}{}^2\bbH(2)&\bbH(4)&\bbC(8)&\bbR(16)\\
2 &\bbR(2)&{}^2\bbR(2)&\bbR(4)&\cellcolor{red!30}\bbC(4)&\bbH(4)&{}^2\bbH(4)&\bbH(8)&\bbC(16)\\
3 &\bbC(2)&\bbR(4)&\cellcolor{red!50}{}^2\bbR(4)&\bbR(8)&\bbC(8)&\bbH(8)&^2\bbH(8)&\bbH(16)\\
4 &\bbH(2)&\cellcolor{red!30}\bbC(4)&\bbR(8)&{}^2\bbR(8)&\bbR(16)&\bbC(16)&\bbH(16)&^2\bbH(16)\\
5 &\cellcolor{red!30}{}^2\bbH(2)&\bbH(4)&\bbC(8)&\bbR(16)&^2\bbR(16)&\bbR(32)&\bbC(32)&\bbH(32)\\
6 &\bbH(4)&{}^2\bbH(4)&\bbH(8)&\bbC(16)&\bbR(32)&{}^2\bbR(32)&\bbR(64)&\bbC(64)\\
7&\bbC(8)&\bbH(8)&{}^2\bbH(8)&\bbH(16)&\bbC(32)&\bbR(64)&{}^2\bbR(64)&\bbR(128)
\end{tabular}
}
%
\end{center}
\small The left superscript denotes direct sum of matrices,
$^2\bbR(2)\equiv\bbR(2)\oplus\bbR(2)$.
\end{block}
\end{column}
\end{columns}
%\cellcolor{blue!25}
\end{frame}
%
%------------ 14 a-------------
\begin{frame}
\vspace{-20pt}
\frametitle<presentation>{Inverses in dimensions $n=p+q\le5$}
\begin{columns}[t]
\begin{column}{1.05\textwidth}
\begin{block}{The case of even multivectors}
\begin{itemize}
\item Even multivectors make subalgebra
\item There is an isomorphisms \begin{equation*}\label{iso}
\cl{p}{q+1}^+\cong\cl{p}{q},\quad\cl{p+1}{q}^+\cong\cl{q}{p}\end{equation*}
\item This can help to get inverse for MV of lower dimensional algebra once
formula of higher dimension algebra is known. 
\item Example for $p+q=4$. Formula of inverse of general MV for $n=4$ is\vskip 1pt 
$\qquad \qquad \qquad \m{A}^{-1}=\frac{\m{A}_{\bar 2,\bar 3}(\m{A}\m{A}_{\bar 2,\bar 3})_{\bar 1,\bar 4}}
{\m{A}\m{A}_{\bar 2,\bar 3}(\m{A}\m{A}_{\bar 2,\bar 3})_{\bar
1,\bar 4}}$
\item If $\m{A}$ consists only of
even grades then $\m{A}$ is invariant to odd grade negations,
so that the above formula for even MV's reduces to\vskip 10pt
$\qquad\qquad\qquad \m{A}^{-1}=\frac{\m{A}_{\bar 2}(\m{A}\m{A}_{\bar 2})_{\bar 4}}
{\m{A}\m{A}_{\bar 2}(\m{A}\m{A}_{\bar 2})_{\bar 4}}$
\end{itemize}
\end{block}
\end{column}
\end{columns}
\end{frame}
% ---- 14 b
\begin{frame}
\vspace{-20pt}
\frametitle<presentation>{Inverses in dimensions $n=p+q\le5$}
\begin{columns}[t]
\begin{column}{1.05\textwidth}
  \begin{block}{The case of even multivectors (continued)}
\begin{itemize}
\item Remember that even grade elements of \cl{1}{3} go to vectors and bivectors of \cl{3}{0} and the pseudoscalar $I_4$ goes to pseudoscalar $I_3$. Thus, in \cl{3}{0}
algebra notation the last formula should be rewritten as\vskip 10pt
$\qquad\qquad\qquad \m{A}^{-1}=\frac{\m{A}_{\bar1,\bar 2}(\m{A}\m{A}_{\bar1,\bar 2})_{\bar 3}}
{\m{A}\m{A}_{\bar 1,\bar 2}(\m{A}\m{A}_{\bar 1,\bar 2})_{\bar
3}}$
\vskip 10pt which coincides with the result for $p+q=3$

\item An alternative formulas for even multivectors exist, which includes even MV multiplication by vector ({\color{red} next slide})

\item The inverse of {\color{blue}even MV} for $p+q=6$ can be written as a self-negated product
\end{itemize}
\end{block}
\end{column}
\end{columns}
\end{frame}


%------------ 15 -------------
\begin{frame}
\vspace{-20pt}
\frametitle<presentation>{Inverses in dimensions $n=p+q\le5$}
\begin{columns}[t]
\begin{column}{1.05\textwidth}
\begin{block}{The case of even multivectors (continued)}
\resizebox{\columnwidth}{!}{%
\begin{tabular}{>{$} l< {$} >{$} l<{$}}
  Cl_{p,q}&\qquad  \bigl(\m{A}^{+}\bigr)^{-1} \\[2pt]\hline\\[-12pt]
p+q=2&\qquad \frac{B \bigl(\m{A}^{+} B\bigr)} {\bigl(\m{A}^{+}
B\bigr)^2}, \qquad B =(\e{i})_{\bar{1}}=-\e{i}
\\[4pt]
  p+q=3 &\qquad \frac{C \bigl(\m{A}^{+} C\bigr)}
{\bigl(\m{A}^{+} C\bigr)^2}, \qquad C = (\m{A}^{+}
\e{i})_{\bar{1}}
\\[4pt]
  p+q=4&\qquad \frac{C \bigl(\m{A}^{+} C\bigr)}
{\bigl(\m{A}^{+} C\bigr)^2}, \qquad  C = (\m{A}^{+}
\e{i})_{\bar{1}}
\\[4pt]
 p+q=5&\qquad \frac{D \bigl(\m{A}^{+} D\bigr)}
{\bigl(\m{A}^{+} D\bigr)^2}, \qquad  D = (\m{A}^{+}
\e{i})_{\bar{3}} \bigl(\m{A}^{+} (\m{A}^{+}
\e{i})_{\bar{3}}\bigr)_{\bar{1}}
\\[4pt]
{\color{red}p+q=6}&\qquad \frac{D \bigl(\m{A}^{+} D\bigr)_{\bar{6}}}
{\bigl(\m{A}^{+} D\bigr)\bigl(\m{A}^{+} D\bigr)_{\bar{6}}}, \quad
D = (\m{A}^{+} \e{i})_{\bar{3}} \bigl(\m{A}^{+} (\m{A}^{+}
\e{i})_{\bar{3}}\bigr)_{\bar{1}}
\end{tabular}
}
\end{block}
\end{column}
\end{columns}
\end{frame}

%------------ 16 -------------
\begin{frame}
\vspace{-20pt}
\frametitle<presentation>{Inverses in dimensions $n=p+q\le5$}
\begin{columns}[t]
\begin{column}{1.05\textwidth}
\begin{block}{The case of even multivectors (continued)}

The denominators in inverse formulas  are the
discriminants of respective multivectors\vskip 5pt
 $s_2=\bigl(\m{A}^{+}
B\bigr)^2/\e{i}^2,\qquad s_4^\prime=\bigl(\m{A}^{+} C\bigr)^2/\e{i}^2,\qquad 
s_4=\bigl(\m{A}^{+} D\bigr)^2/\e{i}^4$
\vskip 5pt and \vskip 5pt $s_8={\bigl(\m{A}^{+}
D\bigr)\bigl(\m{A}^{+} D\bigr)_{\bar{6}}}/\e{i}^4$.\vskip 5pt Here explicit
$\e{i}^2=\pm 1$ and $\e{i}^4=+1$ factors take into account the
correct sign and norm once the orthonormal base vector $\e{i}$ is replaced by arbitrary unnormalized and nonisotropic vector $v$.
\end{block}\vskip 10pt
\begin{block}{Can grade elimination procedure be extended beyond $n=5$?}
  The short answer is 'yes' if we allow {\color{red}multilinear
  combinations} of grade-negated self-products with properly chosen coefficients. 
\end{block}

\end{column}
\end{columns}
\end{frame}
%---------- dim=6 case start
%--------------- 17 ---------
\section[Case $n=p+q=6$;]{Inverse of general MV of algebras of dimension $n=p+q=6$}
\begin{frame}
\frametitle<presentation>{Inverse in dimension $n=p+q=6$}
\begin{columns}[t]
\begin{column}{1.05\textwidth}
\vskip -20pt
\begin{block}{Why single self-negated product does not work in $p+q=6$?}
\begin{itemize}
  \item The culprit appear to be the {\color{red}grade-4} blades. They cannot be eliminated by single sef-product!

\item For $n=6$ we can simultaneously eliminate either grades $1$, $2$, $5$ and $6$ (then $0$, $3$ and $4$ grades survive) or, alternatively,  the grades $2$, $3$ and $6$ (then
grades $0$, $1$, $4$ and $5$ persist).

\item There exists a subalgebra with base  formed by elements $\{1, \e{1256},
\e{1346},\e{2345}\}$ such that any self product of multivector\vskip 10pt
$\qquad \m{B}=a_{0}+a_{47} \e{1256}+a_{49} \e{1346}+a_{52} \e{2345}$\vskip 10pt  by
grade-negated $\m{B}$ will yield new MV having at least one
grade-$4$ base element present.

\item There is {\color{red} only one} such an "uneliminatible" grade part in $p+q=6$ case.
\end{itemize}
\end{block}\vskip -4pt
We will start by trying to eliminate this special grade-4 part.
\end{column}
\end{columns}
\end{frame}



%%%%
%--------------- 18 ---------
\begin{frame}
\frametitle<presentation>{Inverse in dimension $n=p+q=6$}
\begin{columns}[t]
\begin{column}{1.05\textwidth}
\vskip -15pt
\begin{block}{How many terms we need in product?}
\begin{center}\label{TableOfPeriodicity1}
%\setlength{\extrarowheight}{1ex}
%\arraycolsep=0pt
%\parbox[t]{0.9\textwidth}{
%\textit{Period-$8$ table of irreducible
%matrix representations of GA (called the fundamental) for
%Cl$_{p,q}$, where $p$ and  $q$ indicate positive and negative signature.}}
%\setlength{\tabcolsep}{1pt}
%\resizebox{\textwidth}{!}{%
%\renewcommand*{\arraystretch}{1}
%\noalign{\vskip -3pt}\hphantom{m}
\resizebox{\columnwidth}{!}{%  
% Acaus preambule
\begin{tabular}{x{0.1em}| >{$}l<{$} >{$}l<{$} >{$}l<{$} >{$}l<{$}
>{$}l<{$} >{$}l<{$} >{$}l<{$} >{$}l<{$}}
\diag{.1em}{.1em}{\mbox{$\!\!p\,$}}{\mbox{\hphantom{n}\smash{$q$}}}
%Dargio preambule
%\begin{tabular}{>{$} l<{$} | >{$} l<{$} >{$}l<{$} >{$}l<{$} >{$}l<{$}
%>{$}l<{$} >{$}l<{$} >{$}l<{$} >{$}l<{$}}
% tabular body
&0&1&2&3&4&5&6&7\\
\cline{2-9}\\[-12pt]
0 &\bbR&\bbC&\bbH&{}^2\bbH&\bbH(2)&\bbC(4)&\cellcolor{red!50}\bbR(8)&{}^2\bbR(8)\\
1 &{}^2\bbR&\bbR(2)&\bbC(2)&\bbH(2)&{}^2\bbH(2)&\cellcolor{red!30}\bbH(4)&\bbC(8)&\bbR(16)\\
2 &\bbR(2)&{}^2\bbR(2)&\bbR(4)&\bbC(4)&\cellcolor{red!30}\bbH(4)&{}^2\bbH(4)&\bbH(8)&\bbC(16)\\
3 &\bbC(2)&\bbR(4)&{}^2\bbR(4)&\cellcolor{red!50}\bbR(8)&\bbC(8)&\bbH(8)&^2\bbH(8)&\bbH(16)\\
4 &\bbH(2)&\bbC(4)&\cellcolor{red!50}\bbR(8)&{}^2\bbR(8)&\bbR(16)&\bbC(16)&\bbH(16)&^2\bbH(16)\\
5 &{}^2\bbH(2)&\cellcolor{red!30}\bbH(4)&\bbC(8)&\bbR(16)&^2\bbR(16)&\bbR(32)&\bbC(32)&\bbH(32)\\
6 &\cellcolor{red!30}\bbH(4)&{}^2\bbH(4)&\bbH(8)&\bbC(16)&\bbR(32)&{}^2\bbR(32)&\bbR(64)&\bbC(64)\\
7&\bbC(8)&\bbH(8)&{}^2\bbH(8)&\bbH(16)&\bbC(32)&\bbR(64)&{}^2\bbR(64)&\bbR(128)
\end{tabular}
}
%
\end{center}
We need 8 term product. Too many. Fortunately the determinant of matrix representation of scalar+grade-4 multivector is a {\color{red} square!}, i.~e. $\mathrm{det} (\left\langle\m{A}\right\rangle_4)=(\cdots)^2$. Only 4 term product multilinear combination should eliminate grade-4 part.
\end{block}
\end{column}
\end{columns}
\end{frame}
%-------- 19 ------
\begin{frame}
\frametitle<presentation>{Inverse in dimension $n=p+q=6$}
\begin{columns}[t]
\begin{column}{1.06\textwidth}
\begin{block}{Pattern of multilinear combination}
\begin{itemize}
  \item Form prototype of a geometric product of the
form $s_{4+4}=s_{4f} +s_{4g}=$
\begin{equation*}\begin{split}
 &b_1 B f_5\Bigl(f_4(B) f_3\bigl(f_2(B) f_1(B)\bigr)\Bigr)
 +b_2 B g_5\Bigl(g_4(B) g_3\bigl(g_2(B)
  g_1(B)\bigr)\Bigr)\label{s1},
\end{split}\end{equation*}
where $f_j, g_j$ denotes yet undetermined grade-$4$ negations and
$b_1, b_2$ are unknown scalar coefficients of the multilinear
combination. 
\item Not unique. The other three possible choices, however, yield the same answer.
\end{itemize}
\end{block}
Calculate explicit symbolic form of square root of the determinant of matrix
representation of the multivector $\m{B}=a_{0}+a_{47}
\e{1256}+a_{49} \e{1346}+a_{52} \e{2345}$ and compare it with
 all possible negations of grade-$4$ substituted for $f_j$ and $g_j$.
\end{column}
\end{columns}
\end{frame}
%-------- 20 --------
\begin{frame}
\frametitle<presentation>{Inverse in dimension $n=p+q=6$}
\begin{columns}[t]
\begin{column}{1.05\textwidth}
\vskip -20pt
\begin{block}{Modelling the grade negation}
\begin{itemize}
  \item Multiply the grade-4 part by unknown coefficients $p_{4jk}$, where the index $j$ denotes the involution number $f_j,g_j$ in the self-product and $k$ is the
term in multilinear combination (in $p+q=6$ case $k=1$ for $f$ and 
$k=2$ for $g$).
\item The coefficients $p_{ijk}$ acquire the values
$\pm 1$ only, where $-1$ means that involution that changes sign
of grade $i$ is applied and $+1$ means that grade-$i$ remains
unchanged.
\item The answer: $\{b_1=-2/3,\; b_2=-1/3,\;
p_{411}=-1,\; p_{412}=1,\; p_{421}=-1,\; p_{422}=1,\;
p_{431}=-1,\; p_{432}=-1,\; p_{441}=-1,\; p_{442}=1,\;
p_{451}=-1,\; p_{452}=1\}$ {\color{blue} and} $\{b_1=-1/3,\; b_2=-2/3,\;
p_{411}=1,\; p_{412}=-1,\; p_{421}=1,\; p_{422}=-1,\;
p_{431}=-1,\; p_{432}=-1,\; p_{441}=1,\; p_{442}=-1,\;
p_{451}=1,\; p_{452}=-1\}$.
\item Check the answer with the most general {\color{red} scalar + grade-4} multivector.
\end{itemize}
\end{block}
\end{column}
\end{columns}
\end{frame}
%------- 21 --------
\begin{frame}
\frametitle<presentation>{Inverse in dimension $n=p+q=6$}
\begin{columns}[t]
\begin{column}{1.05\textwidth}
\vskip -20pt
\begin{block}{The full discriminant}
\begin{itemize}
  \item Repeat the procedure for involutions that consecutively eliminate all other grades. How to double the number of multivectors?
\item Simultaneusly either grades $1$,
$2$, $5$ and $6$ (then grades $0$, $3$ and $4$ remain), or,
alternatively grades $2$, $3$ and $6$ (then grades $0$, $1$, $4$
and $5$ remain) can be simultaneously eliminated. All in all there
are $\frac{6!}{2!4!}+\frac{6!}{3!3!} +1=36$ elements of grades
$2$, $3$ and $6$ which can be eliminated simultaneously. This is
more than $28$ total elements of grades $1$, $2$, $5$ and $6$ that need to be eliminated. 
\item Therefore, if we replace $\m{B}$ in $s_{4+4}$ pattern by self-negated product $\m{A}\m{A}_{{\bar
2},{\bar 3},{\bar 6}}$ we will have eight required terms. 
\item Then we find the other involutions using the multivector $ a_{0}+a_{1}
\e{1}+a_{22} \e{123}+a_{47} \e{1256}$ and check result with general MV.
\end{itemize}
\end{block}
30 valid meaningfull solutions, which fall into two classes.  Among them there is two solution with minimal number of negated MVs equal to 34.
\end{column}
\end{columns}
\end{frame}
%%%

%------   22  --------
\begin{frame}
\frametitle<presentation>{Inverse in dimension $n=p+q=6$}
\begin{columns}[t]
\begin{column}{1.05\textwidth}
\vskip -20pt
\begin{block}{The answer}
\resizebox{\columnwidth}{!}{%
\begin{tabular}{>{$} l< {$} >{$} l<{$}}
\cl{p}{q}&\qquad  \m{A}^{-1} \\[5pt]\hline\\[-5pt]
p+q=0&\qquad\frac{\m{B}}{\m{A}\m{B}},\qquad\qquad\qquad\qquad\m{B}=1\\[4pt]
%
 p+q=1&\qquad
\frac{\m{B}(\m{A}\m{B})_{\bar{1}}}{\m{A}\m{B}\,(\m{A}\m{B})_{\bar{1}}}
= \frac{\gradeinverse{\m{A}}}{\m{A} \gradeinverse{\m{A}}},\quad\qquad\m{B}=1\\[4pt]
%
  p+q=2&\qquad \frac{\m{C}}{\m{A}\m{C}}=\frac{\cliffordconjugate{\m{A}}}{\m{A} \cliffordconjugate{\m{A}}}
,\qquad\qquad\quad\m{C}=\m{A}_{\bar{1},\bar{2}}\\[4pt]
%
  p+q=3 &\qquad
 \frac{\m{C}(\m{A}\m{C})_{\bar{3}}}{\m{A}\m{C}\,(\m{A}\m{C})_{\bar{3}}}
=\frac{\cliffordconjugate{\m{A}}\gradeinverse{\m{A}}\reverse{\m{A}}}{\m{A}
\cliffordconjugate{\m{A}}\gradeinverse{\m{A}}\reverse{\m{A}}}
,\qquad\m{C}=\m{A}_{\bar{1},\bar{2}}\\[4pt]
%
 p+q=4&\qquad\frac{\m{D}}{\m{A}\m{D}},\qquad\qquad\qquad\qquad
\m{D}=\m{A}_{\bar{2},\bar{3}}
(\m{A}\m{A}_{\bar{2},\bar{3}})_{\bar{1},\bar{4}}\quad \\
   &\qquad\qquad\qquad\qquad\qquad\textrm{or}\quad \m{D}=\m{A}_{\bar{1},\bar{2}}(\m{A}
\m{A}_{\bar{1},\bar{2}})_{\bar{3},\bar{4}}  \\[4pt]
%
 p+q=5&\qquad\frac{\m{D}(\m{A}\m{D})_{\bar{5}}}{\m{A}\m{D}\,
(\m{A}\m{D})_{\bar{5}}},\qquad\qquad\qquad
\m{D}=\m{A}_{\bar{2},\bar{3}}
(\m{A}\m{A}_{\bar{2},\bar{3}})_{\bar{1},\bar{4}}\\[4pt]
%
p+q=6&\cellcolor{red!30}\qquad\frac{\m{G}}{\m{A}\m{G}},\qquad\m{G}=\frac{1}{3} \m{A}_{\bar{2},\bar{3},\bar{6}} \Bigl(\m{H} (\m{H} \m{H})_{\bar{1},\bar{4},\bar{5}} + 2 \bigl(\m{H}_{\bar{4}} (\m{H}_{\bar{4}} \m{H}_{\bar{4}})_{\bar{1},\bar{4},\bar{5}}\bigr)_{\bar{4}}\Bigr)\\
     &\quad\textrm{or}\quad \m{G}=\frac{1}{3} \m{A}_{\bar{2},\bar{3},\bar{6}} \Bigl(\bigl(\m{H} (\m{H} \m{H}_{\bar{1},\bar{5}})_{\bar{4}}\bigr)_{\bar{1},\bar{5}} + 2 \bigl(\m{H}_{\bar{4},\bar{5}} (\m{H}_{\bar{4},\bar{5}} \m{H}_{\bar{1},\bar{4}})_{\bar{4}}\bigr)_{\bar{1},\bar{4}}\Bigr)\\
&\qquad \qquad \qquad\textrm{with}\quad \m{H}=\m{A} \m{A}_{\bar{2},\bar{3},\bar{6}}=\m{A}\reverse{\m{A}}
\end{tabular}
}
\end{block}
\end{column}
\end{columns}
\end{frame}
%%

%---------23 --------
\section[Conclusions;]{Conclusions and perspective}
\begin{frame}
\vspace{-20pt}
%\frametitle<presentation>{Conclusions}
\begin{columns}[t]
\begin{column}{1.06\textwidth}
\begin{block}{Conclusions}
\begin{itemize}\small
\item <1-> {\color{red}Conjecture.} The inverse of a
general MV of arbitrary Clifford algebra can be always expressed
as a multilinear combination of proper number of specially
constructed self-negated products, where in the discriminant expression the initial MV stays
in the left-most or right-most position.
\item <2->  The number of the multivectors in the negated products is determined by dimension of matrix representation of the algebra. 
\item <3-> {\color{red}Hypothesis.} The minimal number of terms in the multilinear combination is related to the number of "uneliminatable" grades in the self product. 
There are no such grades in algebras with $p+q\le 3$.
There is only one grade, namely the grade-$4$, in algebras of
dimension $4$, $5$, $6$ and $7$. The dimensions $4$ and $5$ are exceptional (not enough indices to form cyclic subgroup of the form scalar + grade-4). The algebras with $p+q =8$ and  $p+q =9$ have two such special grades, $4$ and~$8$.
\item <4-> Computational performance? Multiplication of 8 general multivectors requires
$(2^6)^8=281\mspace{2mu}474\mspace{3mu}976\mspace{3mu}710\mspace{4mu}656$
geometric multiplications of base elements in 6-dimensional algebra.
\end{itemize}
\end{block}
\end{column}
\end{columns}
\end{frame}
%-------- 24  ------

\begin{frame}
\vspace{-20pt}
\frametitle<presentation>{\hfill Bibliography\hfill }
\begin{thebibliography}{}

\bibitem{Dadbeh2011}
P.~{Dadbeh},
\newblock {Inverse and determinant in 0 to 5 dimensional {C}lifford algebra},
\newblock arXiv: 1104.0067  (Mar. 2011).

\bibitem{AcusDargys2017}
A.~Acus and A.~Dargys,
\newblock Geometric algebra Mathematica package, 2017,
 \url{https://github.com/ArturasAcus/GeometricAlgebra}.

\bibitem{AcusDargys2017s}
A.~Acus and A.~Dargys,
\newblock submitted to Advances in Applied Clifford algebras
\end{thebibliography}
\vskip 20pt
\begin{center}
\alert{\huge Thank You}
\end{center}
\end{frame}

%\bibliography{Cliffart,Cliffbooks,Cliffrev,graphene}
\end{document}
